\begin{thebibliography}{99}
\bibitem{Arnaiz-Macia-22}V. Arnaiz \& F. Macià - “Concentration of quasimodes for perturbed harmonic oscillators”, 2022 | arXiv | Zbl
\bibitem{Arnaiz-Macia-22-2}V. Arnaiz \& F. Macià - “Localization and delocalization of eigenmodes of harmonic oscillators”, Proc. Amer. Math. Soc. 150 (2022) no. 5, p. 2195-2208 | DOI | Zbl | MR
\bibitem{Anantharaman-08}N. Anantharaman - “Entropy and the localization of eigenfunctions”, Ann. of Math. (2) 168 (2008) no. 2, p. 435-475 | DOI | Zbl | MR
\bibitem{Arnold-61}V. I. Arnol’d - “Some remarks on flows of line elements and frames”, Soviet Math. Dokl. 2 (1961), p. 562-564 | Zbl
\bibitem{Burger-90}M. Burger - “Horocycle flow on geometrically finite surfaces”, Duke Math. J. 61 (1990) no. 3, p. 779-803 | DOI | Zbl | MR
\bibitem{ColinDeVerdiere-79}Y. Colin de Verdière - “Sur le spectre des opérateurs elliptiques à bicaracteristiques toutes périodiques”, Comment. Math. Helv. 54 (1979), p. 508-522 | Zbl | DOI | MR
\bibitem{Colindeverdiere-85}Y. Colin de Verdière - “Ergodicité et fonctions propres du laplacien”, in Séminaire Bony-Sjöstrand-Meyer, 1984–1985, École Polytechnique, Palaiseau, 1985, Exp. No. 13, 8 p. | Zbl | Numdam
\bibitem{Charles-00}L. Charles - Aspects semi-classiques de la quantification géométrique, Ph. D. Thesis, Université Paris Dauphine-PSL, 2000
\bibitem{Landau1}L. Charles - “Landau levels on a compact manifold”, Ann. H. Lebesgue 7 (2024), p. 69-121 | DOI | Numdam | Zbl | MR
\bibitem{Landau2}L. Charles - “On the spectrum of nondegenerate magnetic Laplacians”, Anal. PDE 17 (2024) no. 6, p. 1907-1952 | DOI | Zbl | MR
\bibitem{Charles--23}L. Charles - “Resolvents of Bochner Laplacians in the semiclassical limit”, Comm. Partial Differential Equations 50 (2025) no. 7, p. 899-930 | DOI | Zbl | MR
\bibitem{Cekic-Lefeuvre-24}M. Cekić \& T. Lefeuvre - “Semiclassical analysis on principal bundles” | arXiv | Zbl
\bibitem{Charles-Lefeuvre-25-2}L. Charles \& T. Lefeuvre - “Semiclassical defect measures of magnetic Laplacians on surfaces II. Variable curvature”, 2025, In preparation
\bibitem{Comtet-86}A. Comtet - “On the Landau levels on the hyperbolic plane”, Ann. Physics 173 (1986), p. 185-209 | MR | DOI | Zbl
\bibitem{Dyatlov-Jin-18}S. Dyatlov \& L. Jin - “Semiclassical measures on hyperbolic surfaces have full support”, Acta Math. 220 (2018) no. 2, p. 297-339 | Zbl | MR | DOI
\bibitem{Dyatlov-Jin-Nonnenmacher-22}S. Dyatlov, L. Jin \& S. Nonnenmacher - “Control of eigenfunctions on surfaces of variable curvature”, J. Amer. Math. Soc. 35 (2022) no. 2, p. 361-465 | DOI | MR | Zbl
\bibitem{Dyatlov-22}S. Dyatlov - “Around quantum ergodicity”, Ann. Math. Qué. 46 (2022) no. 1, p. 11-26 | Zbl | DOI | MR
\bibitem{Dyatlov-Zworski-book}S. Dyatlov \& M. Zworski - Mathematical theory of scattering resonances, Grad. Stud. Math., vol. 200, American Mathematical Society, Providence, RI, 2019 | DOI | MR | Zbl
\bibitem{Flaminio-Forni-03}L. Flaminio \& G. Forni - “Invariant distributions and time averages for horocycle flows”, Duke Math. J. 119 (2003) no. 3, p. 465-526 | DOI | MR | Zbl
\bibitem{Furstenberg-73}H. Furstenberg - “The unique ergodicity of the horocycle flow”, in Recent advances in topological dynamics (New Haven, Conn., 1972), Lect. Notes in Math., vol. 318, Springer, Berlin-New York, 1973, p. 95-115 | Zbl | MR | DOI
\bibitem{Ginzburg-96}V. L. Ginzburg - “On the existence and non-existence of closed trajectories for some Hamiltonian flows”, Math. Z. 223 (1996) no. 3, p. 397-409 | DOI | Zbl | MR
\bibitem{Guillemin-Uribe-89}V. Guillemin \& A. Uribe - “Circular symmetry and the trace formula”, Invent. Math. 96 (1989) no. 2, p. 385-423 | DOI | MR | Zbl
\bibitem{Iengo-Li-94}R. Iengo \& D. Li - “Quantum mechanics and quantum Hall effect on Riemann surfaces”, Nuclear Phys. B 413 (1994) no. 3, p. 735-753 | DOI | Zbl | MR
\bibitem{Kordyukov-Taimanov-22}Y. A. Kordyukov \& I. A. Taimanov - “Trace formula for the magnetic Laplacian on a compact hyperbolic surface”, Regul. Chaotic Dyn. 27 (2022) no. 4, p. 460-476 | DOI | MR | Zbl
\bibitem{Lefeuvre-book}T. Lefeuvre - Microlocal analysis in hyperbolic dynamics and geometry, Cours Spécialisés, vol. 32, Société Mathématique de France, Paris, 2026 | MR
\bibitem{Lindenstrauss-06}E. Lindenstrauss - “Invariant measures and arithmetic unique ergodicity. Appendix by E. Lindenstrauss and D. Rudolph”, Ann. of Math. (2) 163 (2006) no. 1, p. 165-219 | DOI | MR | Zbl
\bibitem{Macia-08}F. Macià - “Some remarks on quantum limits on Zoll manifolds”, Comm. Partial Differential Equations 33 (2008) no. 6, p. 1137-1146 | DOI | MR | Zbl
\bibitem{Macia-09}F. Macià - “Semiclassical measures and the Schrödinger flow on Riemannian manifolds”, Nonlinearity 22 (2009) no. 5, p. 1003-1020 | DOI | MR | Zbl
\bibitem{Ma-Ma-23}M. Ma \& Q. Ma - “Semiclassical analysis, geometric representation and quantum ergodicity”, Comm. Math. Phys. 405 (2024) no. 11, article ID 259, 28 pages | DOI | MR
\bibitem{Marklof-Rudnick-00}J. Marklof \& Z. Rudnick - “Quantum unique ergodicity for parabolic maps”, Geom. Funct. Anal. 10 (2000) no. 6, p. 1554-1578 | DOI | MR | Zbl
\bibitem{Macia-Riviere-16}F. Macià \& G. Rivière - “Concentration and non-concentration for the Schrödinger evolution on Zoll manifolds”, Comm. Math. Phys. 345 (2016) no. 3, p. 1019-1054 | DOI | Zbl | MR
\bibitem{Morin-Riviere-24}L. Morin \& G. Rivière - “Quantum unique ergodicity for magnetic Laplacians on T2”, 2024 | arXiv | Zbl
\bibitem{Nonnenmacher-notes}S. Nonnenmacher - “Introduction to semiclassical analysis”, Lecture notes, available at https://www.imo.universite-paris-saclay.fr/ stephane.nonnenmacher/enseign/Cours-Semiclassical-Analysis2019-lecture1-8.pdf, 2019
\bibitem{Paternain-99}G. P. Paternain - Geodesic flows, Progress in Math., vol. 180, Birkhäuser Boston, Inc., Boston, MA, 1999 | DOI | MR | Zbl
\bibitem{Rosenzweig-06}L. Rosenzweig - “Quantum unique ergodicity for maps on the torus”, Ann. Henri Poincaré 7 (2006) no. 3, p. 447-469 | DOI | MR | Zbl
\bibitem{Rudnick-Sarnak-94}Z. Rudnick \& P. Sarnak - “The behaviour of eigenstates of arithmetic hyperbolic manifolds”, Comm. Math. Phys. 161 (1994) no. 1, p. 195-213 | DOI | MR | Zbl
\bibitem{Shnirelman-74-1}A. I. Shnirelman - “Ergodic properties of eigenfunctions.”, Uspehi Mat. Nauk 29 (1974) no. 6(180), p. 181-182 | MR | Zbl
\bibitem{Shnirelman-74-2}A. I. Shnirelman - “Statistical properties of eigenfunctions”, Ann. Math. Qué. 46 (2022) no. 1, p. 3–9 | DOI | MR | Zbl
\bibitem{Tejero-06}C. Tejero Prieto - “Holomorphic spectral geometry of magnetic Schrödinger operators on Riemann surfaces”, Differential Geom. Appl. 24 (2006) no. 3, p. 288-310 | DOI | MR | Zbl
\bibitem{Uribe-Zelditch-93}A. Uribe \& S. Zelditch - “Spectral statistics on Zoll surfaces”, Comm. Math. Phys. 154 (1993) no. 2, p. 313-346 | MR | Zbl | DOI
\bibitem{Weinstein-77}A. Weinstein - “Asymptotics of eigenvalue clusters for the Laplacian plus a potential”, Duke Math. J. 44 (1977), p. 883-892 | DOI | MR | Zbl
\bibitem{Zelditch-87}S. Zelditch - “Uniform distribution of eigenfunctions on compact hyperbolic surfaces”, Duke Math. J. 55 (1987), p. 919-941 | DOI | MR | Zbl
\bibitem{Zelditch-92}S. Zelditch - “On a “quantum chaos” theorem of R. Schrader and M. Taylor”, J. Funct. Anal. 109 (1992) no. 1, p. 1-21 | DOI | MR | Zbl
\bibitem{Zelditch-97}S. Zelditch - “Fine structure of Zoll spectra”, J. Funct. Anal. 143 (1997) no. 2, p. 415-460 | DOI | MR | Zbl
\bibitem{Zelditch-15}S. Zelditch - “Gaussian beams on Zoll manifolds and maximally degenerate Laplacians”, in Spectral theory and partial differential equations (Los Angeles, CA, June 2013), American Mathematical Society, Providence, RI, 2015, p. 169-198 | DOI | MR | Zbl
\bibitem{Zworski-12}M. Zworski - Semiclassical analysis, Grad. Stud. Math., vol. 138, American Mathematical Society, Providence, RI, 2012 | DOI
\end{thebibliography}
